\documentclass[11pt,a4paper]{article}
\usepackage[T2A]{fontenc}
\usepackage[utf8]{inputenc}
\usepackage[english,russian]{babel}
\sloppy
\setlength{\emergencystretch}{2em}
\usepackage[margin=2.1cm]{geometry}
\usepackage{xcolor}
\usepackage{graphicx}
\usepackage{listings}
\usepackage{enumitem}
\usepackage{booktabs}
\usepackage{array}
\newcolumntype{L}[1]{>{\raggedright\arraybackslash}p{#1}}
\usepackage{amsmath}
\usepackage{tikz}
\usetikzlibrary{arrows.meta,positioning,calc,fit,shapes.geometric}
\usepackage{titlesec}
\usepackage{fancyhdr}
\usepackage[hidelinks]{hyperref}
\usepackage{caption}
\captionsetup{font=small,labelfont=bf}
% babel нь \begin{document}-д эдгээрийг дардаг тул AtBeginDocument ашиглана
\AtBeginDocument{%
  \renewcommand{\contentsname}{Гарчиг}%
  \renewcommand{\tablename}{Хүснэгт}%
  \renewcommand{\figurename}{Зураг}%
  \renewcommand{\refname}{Ашигласан материал}%
}

\definecolor{codebg}{RGB}{246,247,249}
\definecolor{accent}{RGB}{0,0,0}

\lstset{
  basicstyle=\ttfamily\footnotesize,
  backgroundcolor=\color{codebg},
  frame=single, rulecolor=\color{gray!40},
  keywordstyle=\color{accent}\bfseries,
  commentstyle=\color{gray}\itshape,
  breaklines=true, columns=fullflexible, keepspaces=true,
  extendedchars=true, inputencoding=utf8,
  literate=
    {А}{{\CYRA}}1
    {Б}{{\CYRB}}1
    {В}{{\CYRV}}1
    {Г}{{\CYRG}}1
    {Д}{{\CYRD}}1
    {Е}{{\CYRE}}1
    {Ё}{{\CYRYO}}1
    {Ж}{{\CYRZH}}1
    {З}{{\CYRZ}}1
    {И}{{\CYRI}}1
    {Й}{{\CYRISHRT}}1
    {К}{{\CYRK}}1
    {Л}{{\CYRL}}1
    {М}{{\CYRM}}1
    {Н}{{\CYRN}}1
    {О}{{\CYRO}}1
    {Ө}{{\CYROTLD}}1
    {П}{{\CYRP}}1
    {Р}{{\CYRR}}1
    {С}{{\CYRS}}1
    {Т}{{\CYRT}}1
    {У}{{\CYRU}}1
    {Ү}{{\CYRY}}1
    {Ф}{{\CYRF}}1
    {Х}{{\CYRH}}1
    {Ц}{{\CYRC}}1
    {Ч}{{\CYRCH}}1
    {Ш}{{\CYRSH}}1
    {Щ}{{\CYRSHCH}}1
    {Ъ}{{\CYRHRDSN}}1
    {Ы}{{\CYRERY}}1
    {Ь}{{\CYRSFTSN}}1
    {Э}{{\CYREREV}}1
    {Ю}{{\CYRYU}}1
    {Я}{{\CYRYA}}1
    {а}{{\cyra}}1
    {б}{{\cyrb}}1
    {в}{{\cyrv}}1
    {г}{{\cyrg}}1
    {д}{{\cyrd}}1
    {е}{{\cyre}}1
    {ё}{{\cyryo}}1
    {ж}{{\cyrzh}}1
    {з}{{\cyrz}}1
    {и}{{\cyri}}1
    {й}{{\cyrishrt}}1
    {к}{{\cyrk}}1
    {л}{{\cyrl}}1
    {м}{{\cyrm}}1
    {н}{{\cyrn}}1
    {о}{{\cyro}}1
    {ө}{{\cyrotld}}1
    {п}{{\cyrp}}1
    {р}{{\cyrr}}1
    {с}{{\cyrs}}1
    {т}{{\cyrt}}1
    {у}{{\cyru}}1
    {ү}{{\cyry}}1
    {ф}{{\cyrf}}1
    {х}{{\cyrh}}1
    {ц}{{\cyrc}}1
    {ч}{{\cyrch}}1
    {ш}{{\cyrsh}}1
    {щ}{{\cyrshch}}1
    {ъ}{{\cyrhrdsn}}1
    {ы}{{\cyrery}}1
    {ь}{{\cyrsftsn}}1
    {э}{{\cyrerev}}1
    {ю}{{\cyryu}}1
    {я}{{\cyrya}}1,
  xleftmargin=0.8em, framexleftmargin=0.8em,
  language=Java, showstringspaces=false
}
\titleformat{\section}{\Large\bfseries\color{accent}}{\thesection}{0.6em}{}
\titleformat{\subsection}{\large\bfseries}{\thesubsection}{0.6em}{}
\pagestyle{fancy}
\fancyhf{}
\fancyfoot[C]{\thepage}
\renewcommand{\headrulewidth}{0.4pt}
\setlength{\headheight}{14pt}
\setlength{\parskip}{3pt}
\fancyhead[L]{\small IV-2-7 $\to$ IV-2-10}
\fancyhead[R]{\small Review --- бөглөсөн}

\begin{document}
\begin{titlepage}
\centering
{\large\bfseries ШИНЖЛЭХ УХААН ТЕХНОЛОГИЙН ИХ СУРГУУЛЬ\par}
\vspace{0.25cm}
{\large\bfseries МЭДЭЭЛЭЛ, ХОЛБООНЫ ТЕХНОЛОГИЙН СУРГУУЛЬ\par}

\vspace{1.4cm}
\includegraphics[width=3.3cm]{../docs/images/shutis_logo.png}

\vspace{1.8cm}
{\Huge\bfseries БИЕ ДААЛТ 1\par}
\vspace{0.25cm}
{\Huge\bfseries ХАРИЛЦАН ШҮҮМЖ\par}

\vspace{1.0cm}
{\large\bfseries Үйлдлийн систем /F.CSM302/\par}

\vfill

\begin{flushleft}
\setlength{\tabcolsep}{0pt}
\begin{tabular}{@{}p{4.3cm}@{}p{10.0cm}@{}}
\textbf{Сэдэв:} &
``Fork-Join Sorting Application. Programming Project 3 (4-р бүлэг),
Recursion tree'' \\[0.85cm]
\textbf{Хуваарилалтын код:} & IV-2-7 \\[0.75cm]
\textbf{Шүүмжилсэн баг:} & IV-2-10 --- The Sleeping Teaching Assistant \\[0.75cm]
\textbf{Хичээл заасан багш:} & Г.Ганбат /F.CS10/ \\[0.45cm]
\textbf{Гүйцэтгэсэн баг:} &
О.Энх /B241940127/ \\
 & Д.Дагиймаа /B200910032/ \\
 & Д.Оюунжаргал /B221940034/ \\
 & Г.Гантогтох /B231930001/ \\
 & Г.Мягмарсүрэн /B241940014/ \\
 & Э.Бүрэнбаяр /B242270084/ \\[0.45cm]
\textbf{Огноо:} & 2026 оны 10 дугаар сарын 26 \\
\end{tabular}
\end{flushleft}

\vspace{1.2cm}
\end{titlepage}

\tableofcontents
\newpage

%============================================================
\section{Ажлын ойлголт}

\textbf{Бодлогын мөн чанар.} Sleeping Teaching Assistant нь сонгодог Sleeping
Barber бодлогын хувилбар. Нэг туслах багш (TA) оффисдоо суудаг: оюутан
байхгүй бол \emph{унтана}, оюутан ирэхэд \emph{сэрч} тусална. TA завгүй үед
ирсэн оюутан коридорын хязгаарлагдмал тооны сандлын нэгэнд сууж хүлээнэ.
Сандал бүгд дүүрэн бол оюутан буцаж явж, хэсэг хугацааны дараа дахин ирнэ.

\textbf{Синхрончлолын асуудал.} Энэ бодлого нь зэрэгцээ ажиллагааны гурван
сонгодог бэрхшээлийг нэг дор агуулдаг:

\begin{enumerate}[nosep,leftmargin=*]
\item \textbf{Алдагдсан сэрээлт (lost wakeup).} Оюутан ``TA унтаж байна уу?''
гэж шалгаад, сэрээх хүртэлх зай завсарт TA өөр оюутныг авч эхэлж болно.
Энэ бол сонгодог \emph{check-then-act} уралдаан.
\item \textbf{Хуваалцсан дарааллын бүрэн бүтэн байдал.} Хүлээлгийн дараалал,
сандлын тоолуурыг олон thread зэрэг өөрчилдөг тул критик муж шаардлагатай.
\item \textbf{Хаалтын дараалал.} Загварыг зогсооход унтаж буй TA болон
хүлээж буй оюутнуудыг бүгдийг сэрээж гаргахгүй бол \texttt{pthread\_join}
мөнхөд гацна.
\end{enumerate}

\textbf{Багийн шийдэл.} Хөдөлгүүрийг \textbf{C / POSIX Pthreads} дээр бичиж,
удирдах цонхыг Python Tkinter дээр хийсэн; хоёулаа \texttt{ctypes}-ээр
холбогдоно. Оюутан бүр болон TA нь \emph{бодит OS thread}. Хуваалцсан бүх
төлөвийг нэг \texttt{pthread\_mutex\_t mu} хамгаална. TA оюутангүй үед
\texttt{work} нэртэй semaphore дээр \emph{бодитоор блоклогдож} унтдаг --- энэ
нь ``унтах''-ыг давталтаар дуурайгаагүй гэдгийг харуулна. Нэмж
\texttt{ta\_tick}, \texttt{ack}, \texttt{started} болон оюутан тус бүрийн
\texttt{ticks[i]}, \texttt{done[i]} semaphore-уудаар нэг \emph{tick} доторх
үйлдлүүдийг тодорхой дараалалд оруулсан. Иймээс ижил seed, ижил тохиргоо
\textbf{байт байтаараа ижил} түүх өгдөг --- бид үүнийг тусад нь шалгасан
(4.2, тест 2).

%============================================================
\section{Туршилтын орчин}

\begin{center}
\small
\begin{tabular}{lL{9.5cm}}
\toprule
Үзүүлэлт & Утга \\
\midrule
OS & Ubuntu 24.04.5 LTS (Linux 6.18) \\
Architecture & x86\_64 \\
Compiler / runtime & gcc 13.3.0, Python 3.13.16, glibc 2.39 \\
Хувилбар & commit \texttt{ec353b0} (\texttt{main}), tag байхгүй \\
Dependency & зөвхөн стандарт libc + pthread; Python pip сан \textbf{шаардлагагүй} \\
Цөмийн тоо & 2 \\
\bottomrule
\end{tabular}
\end{center}

\emph{Тэмдэглэл: энэ орчинд} \texttt{python3-tk} \emph{суугаагүй байсан тул
график интерфейсийг (Queue animation) энэ машин дээр ажиллуулж чадаагүй.
6-р бүлгийн дүрслэлийн үнэлгээг Linux desktop буюу WSLg орчинд тусад нь
хийсэн --- тэр хэсгийг багийн гишүүн биечлэн бөглөнө.}

%============================================================
\section{Build ба Run}

\begin{center}
\small
\begin{tabular}{L{1.7cm}L{5.9cm}L{3.4cm}L{3.4cm}}
\toprule
Алхам & Команд & Хүлээгдсэн үр дүн & Бодит үр дүн \\
\midrule
Clone & \texttt{git clone https://github.com/} \texttt{mendamgalan/IV-2-10-sleeping-ta} &
татагдана & \textbf{ЗӨВ} --- 23 файл, 2.9 MB \\[6pt]
Build & \texttt{make -C sleeping-ta} &
\texttt{libta.so} үүснэ & \textbf{ЗӨВ} --- \texttt{-Wall -Wextra -Wpedantic -Werror} дээр \textbf{нэг ч анхааруулгагүй} \\[6pt]
Test & \texttt{make -C sleeping-ta test} &
бүх тест давна & \textbf{ХЭСЭГЧИЛСЭН} --- 14 хөдөлгүүрийн тест давсан, UI тестийн модуль ачаалагдаагүй (6.1-ийг үзнэ үү) \\[6pt]
Run & \texttt{make -C sleeping-ta run} &
цонх нээгдэнэ & \textbf{ШАЛГААГҮЙ} --- энэ орчинд дэлгэц байхгүй \\
\bottomrule
\end{tabular}
\end{center}

\textbf{Manual-аас гадуур хийх шаардлагатай болсон өөрчлөлт:} \emph{Байхгүй.}
README-ийн алхмууд бичсэнээрээ ажилласан. Энэ нь сайн тал --- доорх 7.1-ийг
үзнэ үү.

%============================================================
\section{Туршилт}

\subsection{Үндсэн тест (тайланд бичсэнийг давтах)}

\begin{center}
\small
\begin{tabular}{L{0.5cm}L{4.3cm}L{4.4cm}L{4.5cm}}
\toprule
\# & Оролт / тохиргоо & Хүлээгдсэн үр дүн & Бодит үр дүн \\
\midrule
1 & \texttt{make test} (өөрийн тестийн багц) &
18 тест, UI-ийнх нь SKIP &
14 давсан, 4 UI тест \textbf{ачаалагдсангүй} (ImportError) \\[6pt]
2 & \texttt{resource\_failure} --- \texttt{pthread\_create} / \texttt{sem\_init} алдаа оруулах &
\texttt{sim\_init} нь $-3$ буцааж цэвэрлэнэ &
\textbf{ЗӨВ} --- 3 нөхцөл бүгд амжилттай \\[6pt]
3 & Сандал $=0$ (\texttt{test\_zero\_chairs}) &
хэн ч хүлээхгүй, шууд буцна &
\textbf{ЗӨВ} --- бидний бие даасан ажиллуулалтаар дараалал үргэлж хоосон, 228 татгалзал, 40 үйлчилгээ \\
\bottomrule
\end{tabular}
\end{center}

\subsection{Нэмэлт тест (өөрсдөө сонгосон)}

\emph{Доорх тестүүдийг багийн тестийн багцаас гадуур, тусдаа бичсэн
\texttt{ctypes} хөтлөгчөөр шууд} \texttt{libta.so} \emph{дээр ажиллуулсан.}

\begin{center}
\small
\begin{tabular}{L{0.5cm}L{3.3cm}L{4.1cm}L{2.7cm}L{3.0cm}}
\toprule
\# & Тохиргоо & Яагаад сонгосон & Хүлээгдсэн & Бодит \\
\midrule
1 & 40 оюутан, 20 сандал, 300 tick & Дээд ачаалал --- дараалал халих, буфер хэтрэх эсэх &
халихгүй & \textbf{ЗӨВ} --- 3813 ирэлт, 3593 татгалзал, snapshot 1340 байт (хязгаар 8192) \\[6pt]
2 & Ижил seed-ээр хоёр удаа & Давтагдах чанарыг баталгаажуулах &
ижил түүх & \textbf{ЗӨВ} --- журнал \textbf{байт байтаараа ижил} \\[6pt]
3 & 20 оюутан, 5 сандал, 500 tick --- 7671 үйл явдлыг шалгах & Инвариант: хүлээгчийн тоо сандлын тооноос хэтрэх ёсгүй &
0 зөрчил & \textbf{ЗӨВ} --- 0 зөрчил \\[6pt]
4 & 6 хүчингүй оролт (оюутан $=0$, $=41$; сандал $=21$; \texttt{amin}$=0$; \texttt{min}$>$\texttt{max}) & Захын ба хүчингүй утгын боловсруулалт &
бүгд татгалзана & \textbf{ЗӨВ} --- 6/6 нь C түвшинд татгалзсан \\[6pt]
5 & \textbf{ThreadSanitizer} --- 3 ажиллуулалт $\times$ 200 tick, зэрэг \texttt{sim\_snapshot()} дуудлагатай & Хуваалцсан төлөв дээр уралдаан (data race) байгаа эсэх &
0 race & \textbf{ЗӨВ} --- \textbf{нэг ч race илрээгүй} \\[6pt]
6 & \textbf{AddressSanitizer + UBSan}, мөн \textbf{Valgrind} \texttt{--leak-check=full} & Санах ойн алдаа, алдагдал &
0 алдаа & \textbf{ЗӨВ} --- алдаа, алдагдал \textbf{байхгүй} \\
\bottomrule
\end{tabular}
\end{center}

%============================================================
\section{Дүрслэлийн үнэлгээ}
\emph{Багшийн оноосон дүрслэл: \textbf{Queue animation}. Доорхийг график
интерфейсийг биечлэн ажиллуулж бөглөв.}

\begin{center}
\small
\begin{tabular}{L{9.4cm}L{2.0cm}L{3.0cm}}
\toprule
Шалгах асуулт & Тийм/Үгүй & Нотолгоо \\
\midrule
Оролт, тохиргоог өөрчлөхөд дүрслэл шинэчлэгдэж байна уу? & & \\[5pt]
Дараалал дахь оюутны тоо бодит төлөвтэй нийцэж байна уу? & & \\[5pt]
TA-ийн төлөв (унтаж / тусалж) зөв харагдаж байна уу? & & \\[5pt]
Тэмдэглэгээ, legend, нэгж ойлгомжтой юу? & & \\[5pt]
Алхамчилж харах эсвэл эргэн шалгах боломжтой юу? & & \\[5pt]
Дүрслэл нь алгоритмыг тайлбарлах \emph{нотолгоо} болж байна уу? & & \\
\bottomrule
\end{tabular}
\end{center}

\textbf{Дүгнэлт:}\\[3mm]
\noindent\rule{\textwidth}{0.4pt}

\vspace{2mm}\noindent\rule{\textwidth}{0.4pt}

%============================================================
\section{Давуу тал}

\textbf{1. TA нь \emph{бодитоор} блоклогддог --- ``унтах''-ыг дуурайгаагүй.}\\[2mm]
\emph{Нотолгоо:} \texttt{src/engine.c:81} --- \texttt{wait\_sem(\&work)}.
Оюутан байхгүй үед TA-ийн thread нь semaphore дээр зогсоно, завсарлагатай
давталтаар (busy-wait) шалгаагүй. Энэ нь сурах бичгийн бодлогын мөн чанарыг
зөв илэрхийлж байна. Мөн \texttt{wait\_sem} нь \texttt{EINTR}-ийг зөв
боловсруулдаг (\texttt{engine.c:40--50}) --- энэ нарийн зүйлийг оюутны ажилд
ховор хардаг.

\vspace{3mm}
\textbf{2. Алдаа оруулах (fault injection) тесттэй.}\\[2mm]
\emph{Нотолгоо:} \texttt{tests/resource\_failure.c} нь макроор
\texttt{pthread\_create} болон \texttt{sem\_init}-ийг орлуулж, 1-р thread,
4-р thread, 6-р semaphore дээр зориудаар алдаа үүсгэнэ. Ингэснээр
\emph{хэсэгчилсэн эхлэлийн цэвэрлэгээ} зөв эсэхийг шалгадаг --- ихэнх багт
огт байдаггүй түвшний тест. Бид ажиллуулахад 3/3 нөхцөл амжилттай.

\vspace{3mm}
\textbf{3. Бүрэн бөгөөд үнэнч баримтжуулалт.}\\[2mm]
\emph{Нотолгоо:} README нь энэ нь \emph{алхамчилсан загвар} бөгөөд ``чөлөөт
OS scheduling-ийн race/stress benchmark \textbf{биш}'' гэдгийг өөрөө ил
хэлсэн. Мөн Reset/хаахаас өмнө CSV экспортлох шаардлагыг анхааруулсан.
Хязгаарлалтаа нуухгүй бичсэн нь ажлын найдвартай байдлыг нэмэгдүүлж байна.
Нэмж GitHub Actions CI нь push бүрд warnings-as-errors build болон Xvfb-тэй
UI тест ажиллуулдаг.

%============================================================
\newpage
\section{Асуудал}

\begin{center}
\small
\begin{tabular}{L{1.9cm}L{5.4cm}L{3.4cm}L{3.5cm}}
\toprule
Priority & Асуудал (давтах алхамтай) & Нөлөө & Сайжруулах санал \\
\midrule
Critical & \emph{Илрээгүй.} & --- & --- \\[8pt]
Major &
\texttt{make test} нь \texttt{python3-tk} суугаагүй орчинд SKIP биш
\textbf{ERROR} болно. Давтах: \texttt{python3-tk}-гүй Linux дээр
\texttt{make -C sleeping-ta test}. Хүлээгдсэн: README-ийн дагуу UI тест
SKIP. Бодит: \texttt{tests/test\_ui.py:3} дахь модулийн түвшний
\texttt{import tkinter} нь \texttt{ImportError} өгч, 14 тест давсан хэдий ч
\texttt{make} нь \texttt{Error 2} буцаана. &
Шүүмжлэгч эсвэл багш ``тест унасан'' гэж буруу дүгнэх эрсдэлтэй.
Дэлгэцгүй сервер дээр CI хийх боломжгүй. &
\texttt{test\_ui.py}-д импортыг хамгаалах:
хамгаалсан \texttt{try / except ImportError} хэлбэрт
оруулж \texttt{unittest.SkipTest} өгөх (8.1-ийг үзнэ үү).
Ингэснээр одоогийн дэлгэцгүй SKIP-тэй адил ажиллана. \\[8pt]
Minor &
\texttt{sim\_recent\_events()} ба \texttt{sim\_events()} нь хуваалцсан
төлөвийг (\texttt{recent[]}, \texttt{recent\_count}, \texttt{journal},
\texttt{journal\_error}) \texttt{mu}-г \textbf{барилгүй} уншина
(\texttt{engine.c:322--363}). Одоогийн хяналтын урсгал зөвхөн tick-ийн
хооронд дууддаг тул аюулгүй --- бид ThreadSanitizer-ээр баталсан. &
Инвариантыг зөвхөн \emph{дуудагчийн сахилга} барьж байна. Хожим UI-г
өөр thread-ээс шинэчилбэл \textbf{чимээгүй} race үүснэ. &
Эдгээр функцийг ч \texttt{mu} дотор оруулах, эсвэл
\texttt{assert(!in\_tick)} маягийн хамгаалалт нэмэх. \\[8pt]
Minor &
\texttt{sim\_snapshot()} дотор \texttt{snprintf}-ийн буцах утгыг шалгалгүй
\texttt{p}-д нэмж байна (\texttt{engine.c:300--316}). \texttt{snprintf} нь
таслагдсан үед \emph{бичих байсан} уртыг буцаадаг тул \texttt{p} нь
буферээс хэтэрч, \texttt{sizeof(out) - (size\_t)p} нь асар том \texttt{size\_t}
болно. &
Одоогийн хязгаараар (40 оюутан $\to$ 1340 байт, буфер 8192) \textbf{хүрэхгүй}
--- бид дээд ачаалал дээр хэмжсэн. Гэхдээ \texttt{MAX\_STUDENTS} нэмэгдвэл
чимээгүй буфер халилт болно. &
Үлдсэн зайг шалгах:
\texttt{if (p >= (int)sizeof(out))} \dots{} эсвэл туслах
\texttt{append()} функц ашиглах. \\[8pt]
Suggestion &
Хувилбарын \textbf{tag байхгүй} (\texttt{git tag} хоосон) ба
\texttt{git shortlog -sn --all} нь \textbf{4} нэр л харуулж байна
(B242270008: 12, Mend-Amgalan: 4, Orgilsakhan: 2, manlaibaatarmaggir: 1). &
``9 дэх долоо хоногт нийтэлсэн хувилбар'' гэдгийг заах боломжгүй. Багийн
гишүүн бүрийн хувь нэмэр \texttt{git log}-оос бүрэн харагдахгүй байж
магадгүй. &
\texttt{v1.0} tag тавих. Commit хийгээгүй гишүүд өөрийн нэрээр дор хаяж нэг
commit хийх, эсвэл хувь нэмрээ README-д тодорхой бичих. \\
\bottomrule
\end{tabular}
\end{center}

\begin{center}
\small
\begin{tabular}{lL{11.0cm}}
\toprule
Түвшин & Утга \\
\midrule
Critical & Үндсэн функц ажиллахгүй эсвэл буруу үр дүн \\
Major & Гол тест, дүрслэл, Manual ноцтой дутуу \\
Minor & Үндсэн ажиллагаанд саад багатай чанарын асуудал \\
Suggestion & Нэмэлт сайжруулалтын санал \\
\bottomrule
\end{tabular}
\end{center}

%============================================================
\section{Сайжруулах санал}

\textbf{1.} \texttt{tests/test\_ui.py}-ийн 3-р мөрөнд байгаа
\texttt{import tkinter as tk} нь \texttt{python3-tk} суугаагүй орчинд бүх
тестийн багцыг унагааж байна. Үүнийг хамгаалсан импорт болгох саналтай:

\begin{lstlisting}[language=Python]
try:
    import tkinter as tk
except ImportError:
    raise unittest.SkipTest('tkinter (python3-tk) суугаагүй')
\end{lstlisting}

Ингэснээр одоогийн ``дэлгэц байхгүй'' SKIP-тэй ижил зан гаргаж, README-д
бичсэн ``UI-ийн 4 тест SKIP болно'' гэсэн амлалт бүх орчинд үнэн болно.

\vspace{3mm}\textbf{2.} \texttt{sim\_snapshot()} дотор \texttt{snprintf}-ийн
үр дүнг хуримтлуулахын өмнө таслагдалтыг шалгах. Одоо хүрэхгүй ч,
\texttt{MAX\_STUDENTS}-ийг 40-өөс дээш болговол чимээгүй буфер халилт болно.
Хамгийн энгийн шийдэл нь үлдсэн зайг шалгадаг жижиг \texttt{append()}
туслах функц бичих.

\vspace{3mm}\textbf{3.} \texttt{v1.0} tag тавьж, README-д багийн гишүүн
бүрийн хариуцсан ажлын хүснэгт нэмэх. Одоо \texttt{git shortlog} 4 нэр
харуулж байгаа тул багш гишүүн бүрийн оролцоог шалгахад хүндрэлтэй.
Хүснэгт нь \texttt{git}-ийн түүхтэй хамт бүрэн зураг өгнө.

%============================================================
\section{Дүгнэлт}

\textbf{1.} \texttt{test\_ui.py}-ийн импортыг хамгаалж, \texttt{make test}
нь \texttt{python3-tk}-гүй орчинд ч бүрэн давдаг болгох. \emph{(Хамгийн
хялбар, хамгийн их нөлөөтэй засвар --- 3 мөр.)}

\vspace{3mm}\textbf{2.} \texttt{sim\_recent\_events()} /
\texttt{sim\_events()}-ийн түгжээний гэрээг кодод илэрхий болгох (түгжээ
дотор оруулах эсвэл assert нэмэх), мөн \texttt{sim\_snapshot()}-ийн
\texttt{snprintf} таслагдалтыг шалгах.

\vspace{3mm}\textbf{3.} \texttt{v1.0} tag тавьж, багийн гишүүдийн хариуцсан
ажлыг README-д баримтжуулах.

\vspace{6mm}
\noindent\textbf{Ерөнхий үнэлгээ:}
Энэ бол \textbf{маш сайн бэлтгэгдсэн} ажил. Хөдөлгүүр нь бодит POSIX thread,
semaphore, mutex ашигласан; TA нь semaphore дээр бодитоор блоклогддог;
хуваалцсан төлөв бүхэлдээ нэг түгжээгээр хамгаалагдсан. Бидний бие даасан
шалгалтаар \textbf{ThreadSanitizer, AddressSanitizer, UBSan, Valgrind
дөрвүүлээ цэвэр}; ижил seed дээр үр дүн байт байтаараа давтагдана;
500 tick, 7671 үйл явдал дотор ``хүлээгч $\le$ сандал'' инвариант
нэг ч удаа зөрчигдөөгүй; хүчингүй оролт 6/6 татгалзагдсан. Компайл нь
\texttt{-Werror} дээр нэг ч анхааруулгагүй.

Илэрсэн цорын ганц ноцтой асуудал нь кодын биш, \textbf{тестийн багцын}
асуудал --- нэмэлт сангүй орчинд SKIP хийх ёстой UI тест нь бүх багцыг
унагааж байна. Үлдсэн хоёр нь одоогоор хүрэхгүй боловч ирээдүйд аюултай
болж болох \emph{далд} эрсдэлүүд. Эдгээрийг засвал ажил бүрэн гүйцэд болно.

\end{document}
